\documentclass[10pt,a4paper]{amsart}
\usepackage[margin=19mm]{geometry}
\usepackage{amsmath,amssymb,mathtools}
\usepackage[scr=rsfs,cal=euler]{mathalfa}
\usepackage{xfrac}
\usepackage[unicode,hidelinks,bookmarks=false]{hyperref}
\newcommand{\ie}{i.e.\ }
\newcommand{\cf}{cf.\ }
\newcommand{\resp}{respectively\ }
\newcommand{\Subsection}{\subsection}
\providecommand{\nobreakdash}{\nobreak\hbox{-}}
\setlength{\emergencystretch}{2em}
\multlinegap0pt
\usepackage{mathtools}
\usepackage{amssymb}
\usepackage{tikz}
\newtheorem{claim}{Claim}
\newtheorem{proposition}{Proposition}
\newcommand\C {{\mathbb C}}
\newcommand{\HH}{{\mathbb H}}
\newcommand\RR{{\rm R}} 
\DeclareMathOperator{\arcosh}{arcosh}
\DeclareMathOperator{\arsinh}{arsinh}
\DeclareMathOperator{\height}{{\rm height}}
\DeclareMathOperator{\id}{id}
\title{Systole, inradius and rigidity of cusped hyperbolic $3$-manifolds}
\author{Stephane Sabourau}
\date{Source: arXiv:2606.06777v1 (CC BY 4.0)}
\begin{document}
\maketitle

\noindent\emph{Source and licence.} Systole, inradius and rigidity of cusped hyperbolic $3$-manifolds. Version arXiv:2606.06777v1 (CC BY 4.0), distributed by arXiv under the Creative Commons Attribution 4.0 International licence, http://creativecommons.org/licenses/by/4.0/, as recorded in the arXiv licence metadata for this exact version and shown on the abstract page https://arxiv.org/abs/2606.06777v1. Full licence notice: \texttt{licenses/arxiv-2606.06777-CC-BY-4.0.txt}.\par\smallskip

\noindent\textit{Editorial scope.} Counted unit: the article's claim claim:1.65 (source0.tex lines 2843--2845). The article's complete original proof of it is delivered with it (source0.tex lines 2848--2979, one \texttt{proof} environment of 132 lines). No mathematics is written, repaired or re-derived: the statement and the proof are the article's own lines, in the article's own notation, and no retained line is substituted or re-flowed.  Retained uncounted as the source-local context the proof needs: lines 506--508; lines 511--513; lines 515--517; lines 533--535; lines 3311--3321; lines 3425--3437.  Nothing was omitted from any retained span, so the package carries no omission record.  No retained line contains a citation, so the wrapper carries no bibliography and no bibliography entry.  No retained line contains a figure environment, an image-inclusion command or drawing code, so no image is delivered and the package declares no asset dependency.  Editorial formatting only, in addition: an AMS \texttt{amsart} preamble that restates the article's own declarations and macros used by the retained lines, namely the environments \texttt{claim}, \texttt{proposition} on separate counters, together with the standard packages those lines need.  The retained environments print in this document as Claim 1, Proposition 1 in order of appearance, whereas the article prints the same items with the numbering of its own full document.  The complete native source of the article as distributed in the arXiv e-print for this exact version is bundled unchanged as \texttt{sources/202212/source0.tex}.
\begin{equation} \label{eq:hyp-eucl} 
d(p,p') = \arcosh \left( 1+ \frac{|p-p'|^2}{2 tt'} \right).
\end{equation}

\begin{equation} \label{eq:hyp-eucl2} 
d(p,p') = \left| \log \left( \frac{t'}{t} \right) \right|.
\end{equation}

\begin{equation} \label{eq:hyp-eucl3} 
d(p,p') = 2 \arsinh \left( \frac{|z-z'|}{2t} \right).
\end{equation}

\begin{equation} \label{eq:h1h2}
| z_1 - z_2 | \geq \sqrt{h_1 h_2}.
\end{equation}

\begin{proposition} \label{prop:diskK}
The Klein bottle~$K(a,b)$ contains an embedded Euclidean open disk of radius~$r$ if and only if
\[
a \geq r, \quad \frac{b}{2} \geq r, \quad a^2 + \left(\frac{b}{2} \right)^2 \geq 4r^2.
\]

Thus, the maximal radius of a Euclidean open disk embedded in~$K(a,b)$ is
\[
\frac{1}{2} \min \left\{ 2a,b,\sqrt{a^2+\left( \frac{b}{2} \right)^2} \right\}.
\]
\end{proposition}

\begin{proposition} \label{prop:hexagon}
Let $\mathcal{H}=ABCDEF$ be a centrally symmetric hexagon in the Euclidean plane with all sides of length~$1$
Assume that the angles at $B,C,E,F$ are equal to $\theta \in [\pi/2,\pi]$. 
Then the maximal distance from a point $P \in \mathcal{H}$ to the set of vertices of~$\mathcal{H}$ is
\[
R=
\begin{cases}
\frac{1}{2\cos \frac{\theta}{2}} & \text{if } \frac{\pi}{2} \leq \theta \leq \frac{2\pi}{3} \\
\sqrt{\frac{1}{4}+\sin^2 \theta} & \text{if } \frac{2\pi}{3} \leq \theta \leq \pi.
\end{cases}
\]
Moreover, for $\theta<\frac{2\pi}{3}$ the maximum is achieved at exactly two symmetric points, whereas for $\theta \geq \frac{2\pi}{3}$ it is achieved uniquely at the center of symmetry~$O$ of $\mathcal{H}$.
\end{proposition}

\begin{claim} \label{claim:1.65}
If $1.65 \leq \theta < \frac{5\pi}{6}$, then $\cosh \RR(M) > \frac{\sqrt{5}}{2}$.
\end{claim}

\begin{proof}
Consider the tiling of the plane by the hexagons in the $\Gamma_\infty$-orbit of~$\mathcal{H}$.
Since we may assume that $\theta \neq \frac{2\pi}{3}$, these hexagons are not regular.
In particular, there is no room for a full size-size horoball to lie inside one of them without intersecting the interior of the full-size horoballs centered at their vertices.

Let $H$ be a horoball (not of full-size) contained in one of these hexagons.
By Proposition~\ref{prop:hexagon}, the center of~$H$ is at distance at most
\[
D=
\begin{cases}
\frac{1}{2 \cos \frac{\theta}{2}} & \text{if } \frac{\pi}{2} \leq \theta \leq \frac{2\pi}{3} \\
\sqrt{\frac{1}{4} + \sin^2 \theta} & \text{if } \frac{2\pi}{3} \leq \theta \leq \frac{5\pi}{6}
\end{cases}
\]
from the center of a full-size horoball.
By the distance--height estimate~\eqref{eq:h1h2}, it follows that
\begin{equation} \label{eq:heightH}
\height(H) \leq D^2.
\end{equation}

Let $p_t$ be the point of height~$t$ above the center~$O$ of the hexagon~$\mathcal{H}$.
We will show that for some~$t>0$ and every $\gamma \in \Gamma \setminus \{ \id \}$,
\[
d(p_t,\gamma.p_t) > 2 \arcosh \left( \frac{\sqrt{5}}{2} \right).
\]

Suppose that $\gamma$ fixes~$H_\infty$, that is, $\gamma \in \Gamma_\infty$.
The point~$p_1$ is the projection of~$p_t$ to~$\partial H_\infty$.
It lies at equal distance from the axes of~$\tau'$ and~$\sigma$.
By Proposition~\ref{prop:diskK}, this implies that
\[
|\gamma.p_t-p_t| = |\gamma.p_1-p_1| \geq 2r
\]
where
\[
r = \frac{1}{2} \min \left\{ 2 \sin \tfrac{\theta}{2}, \sin \tfrac{\theta}{2} \sqrt{1+4 \sin^2 \theta} \right\}.
\]
By the Euclidean--hyperbolic distance formula~\eqref{eq:hyp-eucl3}, this yields
\begin{equation} \label{eq:gamma1}
d(p_t,\gamma.p_t) \geq 2 \arsinh \left( \frac{r}{t} \right) > 2 \arcosh \left( \frac{\sqrt{5}}{2} \right)
\end{equation}
whenever $t<2r$.

\medskip

Suppose that $\gamma$ sends~$H_\infty$ to a horoball~$H$ which is not of full-size.
By the height bound~\eqref{eq:heightH}, the image of~$p_1 \in \partial H_\infty$ under~$\gamma$ has height at most~$D^2$.
Similarly,
\[
\height(\gamma.p_t) \leq \frac{D^2}{t}.
\]
Using the Euclidean--hyperbolic distance formula~\eqref{eq:hyp-eucl2}, we obtain
\begin{equation} \label{eq:gamma2}
d(p_t,\gamma.p_t) \geq \log \left( \frac{t^2}{D^2} \right) > 2 \arcosh \left( \frac{\sqrt{5}}{2} \right)
\end{equation}
whenever $t > \frac{1+\sqrt{5}}{2} \, D$.

\medskip

Before moving on to the next case, observe that 
\[
\frac{1+\sqrt{5}}{2} D < 2r.
\]
Indeed, for $\frac{\pi}{2} \leq \theta \leq \frac{2\pi}{3}$, this is equivalent to
%\[
%\frac{1+\sqrt{5}}{2} \, \frac{1}{2 \cos \frac{\theta}{2}} < 2 \sin \tfrac{\theta}{2} \quad \Longleftrightarrow \quad \tan \tfrac{\theta}{2} > \frac{1+\sqrt{5}}{2}
%\]
\[
\tan \tfrac{\theta}{2} > \frac{1+\sqrt{5}}{2}
\]
which holds since $\tan \frac{\theta}{2} \geq 1$.
For $\frac{2\pi}{3} \leq \theta \leq \frac{5\pi}{6}$, this is equivalent to
%\[
%\frac{1+\sqrt{5}}{2} \, \sqrt{\frac{1}{4} + \sin^2 \theta} < \sin \tfrac{\theta}{2} \, \sqrt{1+4 \sin^2 \theta} \quad \Longleftrightarrow \quad \sin \tfrac{\theta}{2} > \frac{1+\sqrt{5}}{4}.
%\]
\[
\sin \tfrac{\theta}{2} > \frac{1+\sqrt{5}}{4}
\]
which also holds since $\sin \frac{\theta}{2} \geq \frac{\sqrt{3}}{2}$.

Thus, for $t<2r$ sufficiently close to~$2r$, both inequalities~\eqref{eq:gamma1} and~\eqref{eq:gamma2} are satisfied.

\medskip

Suppose that $\gamma$ sends~$H_\infty$ to a full-size horoball~$H$.
Let $d$ denote the Euclidean distance on~$\partial H$ between the projection~$\gamma.p_1$ of~$\gamma.p_t$ to~$\partial H$ and the tangent point of~$H$ and~$H_\infty$.
Then $d \geq R$, otherwise the center of~$\gamma^{-1}.H_\infty$ would be at distance less than~$D$ from~$O$, which is absurd.

By rotational symmetry around the vertical geodesic~$\Delta$ though the center of~$H$, we may assume that $p_t$ and $\gamma.p_t$ lie in the same vertical half-plane bounded by $\Delta$.
More precisely, let $q_1 \in \partial H$ be the point at Euclidean distance~$d$ on~$\partial H$ from the tangency point of~$H$ and~$H_\infty$, measured along~$\partial H$, and lying in the previous half-plane.
Then define~$q_t$ as the point obtained from~$q_1$ by moving a signed hyperbolic distance~$\log(t)$ along the geodesic orthogonal to~$\partial H$ at~$q_1$, directed into~$H$.
It follows that
\[
d(p_t,\gamma_t) \geq d(p_t,q_t).
\]
%where $q_t$ is the unique point of~$\HH^3$ lying in the half-plane delimited by~$\Delta$ containing~$p_t$, whose projection~$q_1$ to~$\partial H$ is at Euclidean distance~$d$ on~$\partial H$ from the tangent point of~$H$ and~$H_\infty$, and whose signed distance (positive inward~$H$) from~$q_1$ is equal to~$\log(t)$.
%\begin{enumerate}
%\item $q_t$ lies in the half-plane delimited by~$\Delta$ containing~$p_t$;
%\item its projection~$q_1$ to~$\partial H$ is at Euclidean distance~$d$ on~$\partial H$ from the tangent point of~$H$ and~$H_\infty$;
%\item the signed distance (positive inward~$H$) between~$q_t$ and~$q_1$ is equal to~$\log(t)$.
%\end{enumerate}
Since we are only seeking a lower bound, we may further assume that the center of~$H$ is as close as possible to~$O$.

To compute~$q_t$ explicitly, we introduce coordinates in which the center of~$H$ is at the origin in~$\C$ and the point~$O$ is represented by~$D$.
In these coordinates, the point~$q_t$ is obtained by applying the inversion with respect to the hemisphere of Euclidean radius~$1$ centered at the origin of~$\C$ to the point~$(d,t)$.
Therefore,
\[
\gamma.p_t = \left( \frac{d}{d^2+t^2}, \frac{t}{d^2+t^2} \right).
\]
Recall that $p_t=(D,t)$ in these coordinates.
By the Euclidean--hyperbolic distance formula~\eqref{eq:hyp-eucl}, we obtain
\[
d(p_t,q_t) = \arcosh \left( 1+\frac{\left( D - \frac{d}{d^2+t^2} \right)^2 + \left( t - \frac{t}{d^2+t^2} \right)^2}{\frac{2 t^2}{d^2+t^2}} \right).
\]
Denote this expression by~$f(D,d,t)$.
One can check that the function $d \mapsto f(D,d,t)$ is increasing for $d \geq D$ and $D^2+t^2 >1$.
Taking $t$ close to~$2r \geq \sqrt{2}$ ensures that $t>1$.
Hence, $d(p_t,q_t) \geq f(D,D,t)$.

\medskip

If $\frac{\pi}{2} \leq \theta \leq \frac{2\pi}{3}$, then $D=\frac{1}{2 \cos \frac{\theta}{2}}$ and $2r = 2 \sin \frac{\theta}{2}$.
The function $\theta \mapsto f(D,D,2r)$ is increasing on this interval and equals $2 \arcosh(\frac{\sqrt{5}}{2})$ for some $\theta_0 \in (1.64,1.65)$.
Thus, for $\theta \geq 1.65$, choosing $t$ sufficiently close to~$2r$ yields
\[
d(p_t,\gamma.p_t) \geq d(p_t,q_t) > 2 \arcosh \left( \frac{\sqrt{5}}{2} \right).
\]

If $\frac{2\pi}{3} \leq \theta < \frac{5\pi}{6}$, then $D=\sqrt{\frac{1}{4} + \sin^2 \theta}$ and $2r = \sin \frac{\theta}{2} \sqrt{1+ 4 \sin^2 \theta}$.
The function $\theta \mapsto f(D,D,2r)$ is decreasing on this interval and equals $2 \arcosh(\frac{\sqrt{5}}{2})$ at $\theta = \frac{5\pi}{6}$.
Thus, for $\theta<\frac{5\pi}{6}$, the same conclusion follows.
\end{proof}

\end{document}
